\documentclass[12pt,a4paper]{article}

% ==================== PACKAGES ====================
\usepackage[utf8]{inputenc}
\usepackage[T1]{fontenc}
\usepackage{lmodern}
\usepackage{geometry}
\geometry{a4paper, margin=1in, top=1.2in, bottom=1.2in}
\usepackage{amsmath,amssymb,amsfonts}
\usepackage{graphicx}
\usepackage{booktabs}
\usepackage{array}
\usepackage{longtable}
\usepackage{xcolor}
\usepackage{listings}
\usepackage{hyperref}
\usepackage{enumitem}
\usepackage{fancyhdr}
\usepackage{titlesec}
\usepackage{tcolorbox}
\usepackage{multicol}
\usepackage{float}

% ==================== COLORS ====================
\definecolor{codegreen}{rgb}{0,0.6,0}
\definecolor{codegray}{rgb}{0.5,0.5,0.5}
\definecolor{codepurple}{rgb}{0.58,0,0.82}
\definecolor{backcolour}{rgb}{0.96,0.96,0.96}
\definecolor{bugred}{rgb}{0.8,0.1,0.1}
\definecolor{fixgreen}{rgb}{0.1,0.6,0.1}
\definecolor{warnorg}{rgb}{0.85,0.55,0.0}
\definecolor{infobox}{rgb}{0.92,0.95,1.0}
\definecolor{warnbox}{rgb}{1.0,0.95,0.88}
\definecolor{dangerbox}{rgb}{1.0,0.92,0.92}

% ==================== CODE LISTINGS ====================
\lstdefinestyle{pythonstyle}{
    backgroundcolor=\color{backcolour},
    commentstyle=\color{codegreen},
    keywordstyle=\color{blue}\bfseries,
    numberstyle=\tiny\color{codegray},
    stringstyle=\color{codepurple},
    basicstyle=\ttfamily\small,
    breakatwhitespace=false,
    breaklines=true,
    captionpos=b,
    keepspaces=true,
    numbers=left,
    numbersep=6pt,
    showspaces=false,
    showstringspaces=false,
    showtabs=false,
    tabsize=4,
    frame=single,
    rulecolor=\color{codegray}
}
\lstset{style=pythonstyle}

% ==================== TCOLORBOX STYLES ====================
\newtcolorbox{infoblock}[1][]{
    colback=infobox, colframe=blue!40!black,
    fonttitle=\bfseries, title=#1,
    boxrule=0.5pt, arc=3pt
}
\newtcolorbox{warnblock}[1][]{
    colback=warnbox, colframe=orange!60!black,
    fonttitle=\bfseries, title=#1,
    boxrule=0.5pt, arc=3pt
}
\newtcolorbox{dangerblock}[1][]{
    colback=dangerbox, colframe=red!60!black,
    fonttitle=\bfseries, title=#1,
    boxrule=0.5pt, arc=3pt
}

% ==================== HEADER/FOOTER ====================
\pagestyle{fancy}
\fancyhf{}
\fancyhead[L]{\small\textit{GlowUpAdvisor --- Comprehensive Technical Report}}
\fancyhead[R]{\small\textit{v1.0 --- August 2026}}
\fancyfoot[C]{\thepage}
\renewcommand{\headrulewidth}{0.4pt}

% ==================== TITLE FORMATTING ====================
\titleformat{\section}{\Large\bfseries\color{blue!60!black}}{}{0em}{\thesection\quad}
\titleformat{\subsection}{\large\bfseries\color{blue!40!black}}{}{0em}{\thesubsection\quad}
\titleformat{\subsubsection}{\normalsize\bfseries}{}{0em}{\thesubsubsection\quad}

% ==================== DOCUMENT ====================
\begin{document}

% ==================== TITLE PAGE ====================
\begin{titlepage}
\centering
\vspace*{2cm}
{\Huge\bfseries GlowUpAdvisor}\\[0.5cm]
{\LARGE AI-Powered Face Shape \& Skin Undertone Analysis\\for Personalized Makeup Recommendations}\\[2cm]
{\Large\itshape Comprehensive Technical Report}\\[0.5cm]
{\large Version 1.0 --- August 2026}\\[3cm]
\rule{\textwidth}{1pt}\\[1cm]
{\large
\begin{tabular}{rl}
\textbf{Project Type:} & R\&D / ML Pipeline in Jupyter Notebook \\
\textbf{Tech Stack:} & Python 3.12 + MediaPipe + OpenCV + scikit-learn + NumPy \\
\textbf{Architecture:} & Still-Photo Analysis (single frame, no AR video) \\
\textbf{Core Output:} & Personalized Makeup Recommendation Card (JSON) \\
\textbf{Product Matching:} & Vectorized CIEDE2000 ($\Delta E_{00}$) \\
\end{tabular}
}\\[2cm]
\rule{\textwidth}{1pt}\\[1cm]
{\small This report consolidates all findings from the Gemini Deep Research phase, the 5-bug technical audit, all mathematical corrections, architectural decisions, and the implementation roadmap.}
\end{titlepage}

\newpage
\tableofcontents
\newpage

% ================================================================
% PART I: PROJECT OVERVIEW
% ================================================================
\part{Project Overview}

\section{What This Project Is}

GlowUpAdvisor is an AI-powered beauty consultation system that takes a single selfie photograph and produces:

\begin{enumerate}[label=\arabic*.]
    \item \textbf{Face Shape Classification} --- Determines whether the user's face is Oval, Round, Square, Heart, Oblong, or Diamond by computing geometric ratios from 478 facial landmark points.
    \item \textbf{Skin Undertone Analysis} --- Extracts skin color from forehead and cheek regions, converts to CIELAB color space, and classifies into one of 12 Personal Color Seasons (e.g., Warm Autumn, Cool Winter).
    \item \textbf{Personalized Makeup Recommendations} --- Generates step-by-step makeup guidance (contour placement, blush angle, lip color, eyeshadow palette, things to avoid) based on the combination of face shape + color season + occasion + skin texture.
    \item \textbf{Real Product Matching} --- Matches recommended shade hex codes to actual cosmetic products (MAC, Romand, 3CE, NARS, etc.) using perceptually accurate color-distance math (CIEDE2000).
\end{enumerate}



% ================================================================
% PART II: TECH STACK
% ================================================================
\part{Technology Stack}

\section{Chosen Technologies}

\begin{longtable}{@{}l l p{8cm}@{}}
\toprule
\textbf{Layer} & \textbf{Technology} & \textbf{Rationale} \\
\midrule
\endhead
Runtime & Python 3.12 + Jupyter & User preference; cell-by-cell iteration; inline visualization \\
Face Landmarks & MediaPipe Face Mesh & 478 3D points; free; local inference; no API key; Google-backed \\
Image Processing & OpenCV (\texttt{cv2}) & Industry standard; fast BGR/RGB/Lab conversion; ROI operations \\
Color Science & NumPy (custom) & Replaces slow \texttt{colormath}; 168$\times$ faster at $N{=}1000$ products \\
ML Classification & scikit-learn & Random Forest / SVM / XGBoost for face shape \& season (Phase 3) \\
Data Handling & Pandas + NumPy & Product DB as DataFrame; skin metrics as arrays \\
Visualization & Matplotlib + Seaborn & ROC curves, confusion matrices, color palettes, face annotations \\
Interactive Demo & Gradio or ipywidgets & Quick web UI from notebook (Phase 4) \\
\bottomrule
\end{longtable}

\subsection{Why Not These Alternatives?}

\begin{itemize}
    \item \textbf{Why not TensorFlow/PyTorch?} Overkill for Phase 1. MediaPipe already runs a neural face mesh model internally. Custom deep models are only needed if rule-based/sklearn classifiers underperform on the season classification task.
    \item \textbf{Why not \texttt{colormath}?} Pure Python loops. At $N{=}1000$ products, \texttt{delta\_e\_cie2000()} takes 185 ms per query. Our vectorized NumPy implementation does it in 1.1 ms. See Bug \#3 in Section~\ref{sec:audit}.
    \item \textbf{Why not a web stack (React + Vite)?} Valid for production deployment (Phase 5). For R\&D, Jupyter allows faster experimentation, inline plots, and easier debugging.
    \item \textbf{Why not dlib?} MediaPipe gives 478 landmarks vs. dlib's 68. The extra density is critical for accurate jaw/cheekbone/forehead ratio computation.
\end{itemize}


% ================================================================
% PART III: SYSTEM ARCHITECTURE
% ================================================================
\part{System Architecture \& Pipeline}

\section{End-to-End Pipeline Overview}

The system is a 6-stage sequential pipeline. Each stage is implemented as one or more Jupyter Notebook cells.

\begin{center}
\fbox{\parbox{0.9\textwidth}{
\centering
\textbf{Input Photo (sRGB)} \\
$\downarrow$ \\
\textbf{Stage 1:} White Balance Correction (Gray World) \\
$\downarrow$ \\
\textbf{Stage 2:} MediaPipe Face Mesh (478 3D landmarks) \\
$\downarrow$ \hspace{3cm} $\downarrow$ \\
\textbf{Stage 3:} Face Shape Classification \hspace{1cm} \textbf{Stage 4:} Skin ROI + CIELAB \\
$\downarrow$ \hspace{3cm} $\downarrow$ \\
\hspace{3cm} \textbf{Stage 5:} ITA\textdegree{} + Season Classification \\
$\downarrow$ \hspace{3cm} $\downarrow$ \\
\textbf{Stage 6:} GlowUpAdvisor Engine + Product Matching \\
$\downarrow$ \\
\textbf{Output:} Personalized Makeup Card (JSON) + Matched SKUs
}}
\end{center}

\section{Stage 1: Image Preprocessing \& White Balance}

\subsection{The Problem}
User selfies are captured under wildly varying ambient lighting---warm tungsten ($\sim$2700K), cool fluorescent ($\sim$4000K), daylight ($\sim$5500K), mixed. Without correction, $L^*a^*b^*$ values shift by up to $\Delta E_{00} \approx 8$--$12$, making skin tone classification unreliable.

\subsection{Gray World White Balance}
Assumes the average color in any natural scene should be neutral gray. Scales R and B channels:

\begin{equation}
R' = R \times \frac{\bar{G}}{\bar{R}}, \qquad B' = B \times \frac{\bar{G}}{\bar{B}}
\label{eq:grayworld}
\end{equation}

where $\bar{R}, \bar{G}, \bar{B}$ are the mean pixel values across the image.

\begin{lstlisting}[language=Python, caption={Gray World White Balance}]
import cv2
import numpy as np

def gray_world_wb(img_bgr):
    """Apply Gray World white balance correction."""
    b, g, r = cv2.split(img_bgr.astype(np.float32))
    avg_b, avg_g, avg_r = b.mean(), g.mean(), r.mean()
    r_corrected = np.clip(r * (avg_g / (avg_r + 1e-6)), 0, 255)
    b_corrected = np.clip(b * (avg_g / (avg_b + 1e-6)), 0, 255)
    return cv2.merge([b_corrected, g, r_corrected]).astype(np.uint8)
\end{lstlisting}

\begin{warnblock}[Limitation]
Gray World assumes the scene contains a diverse distribution of colors. A close-up selfie against a colored wall may violate this assumption. Alternative: ask users to include a gray reference card, or use a face-skin-only ROI for white balance estimation.
\end{warnblock}


\section{Stage 2: Face Mesh Extraction (MediaPipe)}

\subsection{What MediaPipe Returns}
MediaPipe Face Mesh detects 478 3D facial landmarks per face, with coordinates normalized to $[0, 1]$ relative to the image dimensions:
\begin{itemize}
    \item \texttt{landmark.x} --- horizontal position (0 = left edge, 1 = right edge)
    \item \texttt{landmark.y} --- vertical position (0 = top edge, 1 = bottom edge)
    \item \texttt{landmark.z} --- depth estimate (relative to face center, negative = closer to camera)
\end{itemize}

\subsection{Critical Landmark Indices}

\begin{table}[H]
\centering
\caption{Key MediaPipe Landmark Indices Used in GlowUpAdvisor}
\begin{tabular}{@{}l l p{6cm}@{}}
\toprule
\textbf{Facial Region} & \textbf{Indices} & \textbf{Usage} \\
\midrule
Top of forehead (hairline) & 10 & Face length reference (top) \\
Chin (bottom) & 152 & Face length reference (bottom) \\
Left cheekbone & 234 & Face width (cheekbone level) \\
Right cheekbone & 454 & Face width (cheekbone level) \\
Left jawline & 172 & Jaw width measurement \\
Right jawline & 397 & Jaw width measurement \\
Left forehead & 338 & Forehead width \\
Right forehead & 297 & Forehead width \\
Forehead center (skin ROI) & 10, 151 & Color sampling zone \\
Left cheek (skin ROI) & 50 & Color sampling zone \\
Right cheek (skin ROI) & 280 & Color sampling zone \\
Outer lip contour & 61--291 & Lip shape + virtual lipstick mask \\
Inner lip contour & 78--308 & Inner lip boundary \\
\bottomrule
\end{tabular}
\end{table}

\begin{lstlisting}[language=Python, caption={MediaPipe Face Mesh Extraction}]
import mediapipe as mp
import cv2

def extract_face_mesh(img_rgb):
    """Extract 478 facial landmarks from an RGB image."""
    mp_face_mesh = mp.solutions.face_mesh
    with mp_face_mesh.FaceMesh(
        static_image_mode=True,
        max_num_faces=1,
        refine_landmarks=True,  # enables 478 (vs 468) landmarks
        min_detection_confidence=0.5
    ) as face_mesh:
        results = face_mesh.process(img_rgb)
        
    if not results.multi_face_landmarks:
        raise ValueError("No face detected in image.")
    
    landmarks = results.multi_face_landmarks[0]
    h, w = img_rgb.shape[:2]
    
    # Convert normalized coords to pixel coords
    points = np.array([
        (lm.x * w, lm.y * h, lm.z * w)  # z scaled by width
        for lm in landmarks.landmark
    ])
    return points  # shape: (478, 3)
\end{lstlisting}


\section{Stage 3: Face Shape Classification}

\subsection{Geometric Ratios}

Three ratios computed from landmark Euclidean distances:

\begin{align}
r_1 &= \frac{d(\text{forehead top}, \text{chin})}{d(\text{left cheekbone}, \text{right cheekbone})} = \frac{d(10, 152)}{d(234, 454)} \label{eq:r1} \\[6pt]
r_2 &= \frac{d(\text{left jaw}, \text{right jaw})}{d(\text{left cheekbone}, \text{right cheekbone})} = \frac{d(172, 397)}{d(234, 454)} \label{eq:r2} \\[6pt]
r_3 &= \frac{d(\text{left forehead}, \text{right forehead})}{d(\text{left jaw}, \text{right jaw})} = \frac{d(338, 297)}{d(172, 397)} \label{eq:r3}
\end{align}

where $d(i, j) = \sqrt{(x_i - x_j)^2 + (y_i - y_j)^2}$ (2D Euclidean distance, ignoring $z$).

\subsection{Classification Rules (v1 --- Rule-Based)}

\begin{table}[H]
\centering
\caption{Face Shape Decision Rules from Geometric Ratios}
\begin{tabular}{@{}l l l l p{4cm}@{}}
\toprule
\textbf{Shape} & $r_1$ (Length/Width) & $r_2$ (Jaw/Cheek) & $r_3$ (Forehead/Jaw) & \textbf{Visual Description} \\
\midrule
Oval & $1.1$--$1.3$ & $0.75$--$0.85$ & $\sim 1.0$ & Balanced proportions, slightly longer than wide \\
Round & $<1.1$ & $0.85$--$0.95$ & $\sim 1.0$ & Nearly equal length and width, soft jaw \\
Square & $<1.15$ & $>0.90$ & $\sim 1.0$ & Strong angular jaw, wide and short \\
Heart & $1.05$--$1.25$ & $<0.75$ & $>1.15$ & Wide forehead, narrow chin \\
Oblong & $>1.3$ & $0.80$--$0.90$ & $\sim 1.0$ & Significantly longer than wide \\
Diamond & $1.1$--$1.25$ & $<0.80$ & $0.85$--$1.0$ & Narrow forehead \& jaw, wide cheekbones \\
\bottomrule
\end{tabular}
\label{tab:face_shape_rules}
\end{table}

\begin{warnblock}[Known Limitation]
These thresholds are heuristic approximations. They have \textbf{not been validated} on a labeled dataset. Phase 3 replaces them with a trained ML classifier (Random Forest or SVM on a dataset of $N \geq 500$ labeled face images).
\end{warnblock}


\section{Stage 4: Skin Color Extraction \& CIELAB Conversion}

\subsection{ROI Extraction}
Three skin regions are sampled (forehead center, left cheek, right cheek) using circular patches centered on the corresponding landmarks. These regions are chosen because they:
\begin{itemize}
    \item Are relatively flat (minimal curvature $\rightarrow$ less shading artifact)
    \item Are far from lips, eyebrows, and eyes (avoids cosmetic contamination)
    \item Represent different blood perfusion zones (forehead = more melanin, cheeks = more hemoglobin)
\end{itemize}

The median $L^*a^*b^*$ across all three patches is used as the final skin color measurement (median is robust to outlier pixels from moles, blemishes, or specular highlights).

\subsection{Color Space Conversion Pipeline}

\begin{center}
\fbox{\textbf{sRGB} $\xrightarrow{\text{gamma}}$ \textbf{Linear RGB} $\xrightarrow{M_{sRGB}}$ \textbf{CIEXYZ (D65)} $\xrightarrow{f(t)}$ \textbf{CIELAB}}
\end{center}

\textbf{Step 1: sRGB $\to$ Linear RGB} (inverse gamma correction):
\begin{equation}
C_{\text{linear}} = \begin{cases}
\frac{C_{\text{sRGB}}}{12.92} & \text{if } C_{\text{sRGB}} \leq 0.04045 \\[4pt]
\left(\frac{C_{\text{sRGB}} + 0.055}{1.055}\right)^{2.4} & \text{otherwise}
\end{cases}
\end{equation}

\textbf{Step 2: Linear RGB $\to$ CIEXYZ (D65 white point)}:
\begin{equation}
\begin{bmatrix} X \\ Y \\ Z \end{bmatrix} =
\begin{bmatrix}
0.4124 & 0.3576 & 0.1805 \\
0.2126 & 0.7152 & 0.0722 \\
0.0193 & 0.1192 & 0.9505
\end{bmatrix}
\begin{bmatrix} R_{\text{lin}} \\ G_{\text{lin}} \\ B_{\text{lin}} \end{bmatrix}
\end{equation}

\textbf{Step 3: CIEXYZ $\to$ CIELAB}:
\begin{align}
L^* &= 116 \cdot f\!\left(\frac{Y}{Y_n}\right) - 16 \\
a^* &= 500 \cdot \left[f\!\left(\frac{X}{X_n}\right) - f\!\left(\frac{Y}{Y_n}\right)\right] \\
b^* &= 200 \cdot \left[f\!\left(\frac{Y}{Y_n}\right) - f\!\left(\frac{Z}{Z_n}\right)\right]
\end{align}
where $f(t) = t^{1/3}$ if $t > \delta^3$ (with $\delta = 6/29$), else $f(t) = \frac{t}{3\delta^2} + \frac{4}{29}$.

D65 white point: $(X_n, Y_n, Z_n) = (95.047, 100.0, 108.883)$.

\begin{infoblock}[Practical Note]
OpenCV's \texttt{cv2.cvtColor(img, cv2.COLOR\_BGR2Lab)} handles this entire pipeline internally. However, understanding each step is essential for debugging color mismatches and for implementing the Bradford CAT (Section~\ref{sec:bradford}).
\end{infoblock}


\section{Stage 5: Skin Metrics \& Season Classification}

\subsection{Individual Typology Angle (ITA\textdegree{})}
\label{sec:ita}

The ITA\textdegree{} quantifies overall skin lightness and pigmentation level on a single angular scale:

\begin{equation}
\boxed{ITA^{\circ} = \arctan\!\left(\frac{L^* - 50}{b^*}\right) \times \frac{180}{\pi}}
\label{eq:ita_correct}
\end{equation}

\begin{dangerblock}[CRITICAL: Use atan, NOT atan2]
The formula uses the \textbf{single-argument} $\arctan(y/x)$, NOT the two-argument $\text{atan2}(y, x)$. Using \texttt{atan2} introduces a $180^{\circ}$ quadrant flip when $b^*$ is negative (cool/bluish skin), causing Winter skin tones to be misclassified as Spring. This was Bug \#1 in the audit (Section~\ref{sec:bug1}).
\end{dangerblock}

\begin{table}[H]
\centering
\caption{ITA\textdegree{} Classification Ranges (Chardon et al., 1991)}
\begin{tabular}{@{}l c l@{}}
\toprule
\textbf{Category} & \textbf{ITA\textdegree{} Range} & \textbf{Approximate Fitzpatrick} \\
\midrule
Very Light & $> 55^{\circ}$ & Type I \\
Light & $41^{\circ}$--$55^{\circ}$ & Type II \\
Intermediate & $28^{\circ}$--$41^{\circ}$ & Type III \\
Tan & $10^{\circ}$--$28^{\circ}$ & Type IV \\
Brown & $-30^{\circ}$--$10^{\circ}$ & Type V \\
Dark & $< -30^{\circ}$ & Type VI \\
\bottomrule
\end{tabular}
\end{table}

\subsection{Additional Skin Metrics}

\begin{itemize}
    \item \textbf{$a^*$ value} (red--green axis): Higher $a^* \to$ more hemoglobin / surface redness. Warm undertones tend to have lower $a^*$; cool undertones higher.
    \item \textbf{$b^*$ value} (yellow--blue axis): Higher $b^* \to$ more melanin / yellow warmth. Warm $\to$ high $b^*$; Cool $\to$ low $b^*$.
    \item \textbf{Melanin Index}: $M = 100 \times \log_{10}(1 / R_{\text{red}})$ --- proxy for melanin concentration from the red channel reflectance.
    \item \textbf{Hemoglobin Index}: $H = 100 \times \log_{10}(R_{\text{red}} / R_{\text{green}})$ --- proxy for blood perfusion.
\end{itemize}

\subsection{12-Season Personal Color Classification}

The 12-season system expands the traditional 4-season (Spring/Summer/Autumn/Winter) by adding depth and clarity modifiers:

\begin{table}[H]
\centering
\caption{12 Personal Color Seasons --- Mapping from Skin Metrics}
\begin{tabular}{@{}l l l l@{}}
\toprule
\textbf{Season} & \textbf{Warmth} & \textbf{Depth} & \textbf{Key Discriminator} \\
\midrule
Light Spring & Warm & Light & High $b^*$, high $L^*$, high ITA\textdegree{} \\
Warm Spring & Warm & Medium & High $b^*$, moderate $L^*$ \\
Bright Spring & Neutral-Warm & Light & High chroma, high $L^*$ \\
Light Summer & Cool & Light & Low $b^*$, high $L^*$, high $a^*$ \\
Cool Summer & Cool & Medium & Low $b^*$, muted chroma \\
Soft Summer & Neutral-Cool & Medium & Low chroma, grayed tones \\
Warm Autumn & Warm & Deep & High $b^*$, low $L^*$, low ITA\textdegree{} \\
Deep Autumn & Warm & Deep-Rich & Very high $b^*$, very low $L^*$ \\
Soft Autumn & Neutral-Warm & Medium & Moderate $b^*$, muted \\
Bright Winter & Neutral-Cool & Light & High chroma, cool-leaning, high $L^*$ \\
Cool Winter & Cool & Deep & Very low $b^*$, low $L^*$, high $a^*$ \\
Deep Winter & Cool & Deep-Rich & Very low $b^*$, very low $L^*$ \\
\bottomrule
\end{tabular}
\end{table}

\begin{warnblock}[Priority \#1 Technical Risk]
The 12-season rule-based thresholds are the \textbf{single most fragile component} of the system. Wrong season $\to$ every downstream recommendation is wrong (lip color, eyeshadow palette, blush tone). These thresholds \textbf{must} be empirically validated via the ROC/AUC calibration framework (Section~\ref{sec:calibration}).
\end{warnblock}


\section{Stage 6: GlowUpAdvisor Engine}

\subsection{Advice Matrix Structure}

The core knowledge base maps \texttt{(face\_shape, season)} $\to$ makeup techniques:

\begin{itemize}
    \item \textbf{Base layer}: $6 \times 12 = 72$ entries with contour, blush, lip, eyes, brow, avoid list
    \item \textbf{Occasion modifier}: 3 tiers (daily / office / glam) $\to$ intensity, lip finish, eye depth
    \item \textbf{Texture modifier}: 3 types (oily / dry / combination) $\to$ skin prep, setting method
    \item \textbf{Total unique profiles}: $6 \times 12 \times 3 \times 3 = 648$
\end{itemize}

Rather than hardcoding 648 entries, the engine uses a \textbf{Strategy Pattern} with composable modifiers merged at runtime.

\subsection{Collision-Safe Dictionary Merge}

\begin{lstlisting}[language=Python, caption={Strategy Pattern with Key Collision Detection}]
def resolve_glowup_profile(shape, season, occasion="daily", texture="combination"):
    base = ADVICE_MATRIX[(shape, season)]
    occ  = OCCASION_MODIFIERS[occasion]
    tex  = TEXTURE_MODIFIERS[texture]
    
    # CRITICAL: detect key collisions before merge
    collisions = occ.keys() & tex.keys()
    assert not collisions, f"Modifier key collision: {collisions}"
    
    return {
        "profile": {"face_shape": shape, "season": season,
                     "occasion": occasion, "texture": texture},
        "techniques": base,
        "execution_parameters": {**occ, **tex}
    }
\end{lstlisting}

\begin{dangerblock}[Bug \#5 --- Silent Key Collision]
Without the collision check, \texttt{\{**occ, **tex\}} silently drops duplicate keys (the second dict wins). If both \texttt{OCCASION\_MODIFIERS} and \texttt{TEXTURE\_MODIFIERS} define a key like \texttt{"finish"}, the occasion value is silently overwritten by the texture value with \textbf{zero error or warning}.
\end{dangerblock}


\section{Resolved System Optimization Decisions}
\label{sec:resolved_decisions}

To maximize optical precision, user experience, and long-term architectural maintainability, four core technical trade-offs have been resolved:

\begin{table}[H]
\centering
\caption{Resolved Technical Decisions \& System Optimizations}
\begin{tabular}{@{}p{4.5cm} p{5.5cm} p{4.5cm}@{}}
\toprule
\textbf{Technical Decision} & \textbf{Selected Optimal Strategy} & \textbf{System Benefit} \\
\midrule
1. Season Classifier Architecture & Modular Interface (\texttt{SeasonClassifier} base class) & Phase 1 ships with rule-based heuristics in \texttt{config.py}; Phase 3 drops in ML model with 0 code refactoring. \\
2. White Balance Strategy & Skin-ROI Masked White Balance & Restricts chromatic adaptation to facial skin landmarks, eliminating background color distortion. \\
3. Face Shape Categorization & Softmax Probability Distribution ($P(\text{shape} \mid r_1, r_2, r_3)$) & Replaces rigid ratio boundaries with continuous scores, preventing advice flips on edge-case faces. \\
4. SKU Match Quality & Top-3 Candidates with $\Delta E_{00}$ Confidence Tiers & Provides user options categorized into Exact Match ($<1.0$), Great Match ($<2.0$), and Close Alternative ($<3.5$). \\
\bottomrule
\end{tabular}
\end{table}

\subsection{Skin-ROI Masked White Balance}
Standard Gray World white balance computes mean channel illuminants $(\bar{R}, \bar{G}, \bar{B})$ over the entire image frame. In real-world selfies, colored backgrounds (e.g., blue walls, indoor lighting fixtures, clothing) distort the estimated illuminant vector.

The optimized pipeline extracts facial skin landmark polygons $\mathcal{P}_{\text{skin}}$ (forehead and cheeks) via MediaPipe and evaluates white balance strictly over skin pixels:
\begin{equation}
\bar{R}_{\text{skin}} = \frac{1}{|\mathcal{P}_{\text{skin}}|} \sum_{(x,y) \in \mathcal{P}_{\text{skin}}} R(x,y), \qquad \bar{B}_{\text{skin}} = \frac{1}{|\mathcal{P}_{\text{skin}}|} \sum_{(x,y) \in \mathcal{P}_{\text{skin}}} B(x,y)
\end{equation}

\subsection{Fuzzy Face Shape Probability Model}
To prevent jarring advice flips when geometric ratios fall near threshold boundaries (e.g., $r_1 = 1.14$ between Round and Oval), the system evaluates a normalized Gaussian membership function for each candidate shape $s$:
\begin{equation}
P(s \mid r_1, r_2, r_3) = \frac{\exp\left(-\frac{1}{2} \sum_{i=1}^3 \left(\frac{r_i - \mu_{s,i}}{\sigma_{s,i}}\right)^2\right)}{\sum_{s'} \exp\left(-\frac{1}{2} \sum_{i=1}^3 \left(\frac{r_i - \mu_{s',i}}{\sigma_{s',i}}\right)^2\right)}
\end{equation}
The engine outputs a primary recommendation based on $\arg\max_s P(s)$ alongside secondary soft recommendations.

% ================================================================
% PART IV: TECHNICAL AUDIT
% ================================================================
\part{Technical Audit --- 5 Critical Bugs}
\label{sec:audit}

This section documents every bug found during the audit, its root cause, severity, mathematical proof of the error, and the fix applied.

\section{Bug \#1: \texttt{atan2} vs \texttt{atan} in ITA\textdegree{} (Critical)}
\label{sec:bug1}

\begin{dangerblock}[Severity: CRITICAL --- Silent Misclassification]
This bug causes cool-toned skin (Winter/Summer) to be classified as warm-toned (Spring/Autumn). Every downstream recommendation is wrong.
\end{dangerblock}

\subsection{The Error}
The original code used Python's \texttt{math.atan2(y, x)}:
\begin{lstlisting}[language=Python]
# WRONG --- introduces quadrant error on cool skin tones
ita = math.degrees(math.atan2((l_val - 50.0), b_val + 1e-5))
\end{lstlisting}

\subsection{Why It's Wrong}
\texttt{atan2(y, x)} returns angles in $[-\pi, \pi]$ across all four quadrants. When $b^* < 0$ (characteristic of cool/bluish skin):
\begin{itemize}
    \item The point $(b^*, L^*{-}50)$ falls in Quadrant II or III of the Cartesian plane
    \item \texttt{atan2} returns an angle offset by $\pm 180^{\circ}$ from the correct ITA\textdegree{}
    \item A Winter skin with true $ITA^{\circ} \approx 20^{\circ}$ gets computed as $ITA^{\circ} \approx -160^{\circ}$
\end{itemize}

\subsection{The Fix}
\begin{lstlisting}[language=Python]
# CORRECT --- single-argument arctan per ITA definition
ita = math.degrees(math.atan((l_val - 50.0) / (b_val + 1e-5)))
\end{lstlisting}


\section{Bug \#2: No Temporal Landmark Stabilization (Critical for AR)}

\subsection{The Error}
Per-frame independent landmark detection $\to$ coordinates jitter $\pm 3$--$5$ pixels $\to$ AR overlay flickers visually.

\subsection{The Fix (Designed)}
Exponential Moving Average (EMA):
\begin{equation}
s_t = \alpha \cdot x_t + (1 - \alpha) \cdot s_{t-1}, \qquad \alpha \approx 0.20\text{--}0.30 \text{ for 30 FPS}
\end{equation}

\subsection{Actual Resolution}
\textbf{Architectural pivot to still-photo analysis.} Temporal jitter is fundamentally impossible in a single-frame pipeline. Bug \#2 is resolved by design, not by code.


\section{Bug \#3: \texttt{colormath} Pure-Python Bottleneck (High)}

\subsection{The Error}
\texttt{delta\_e\_cie2000()} from the \texttt{colormath} library iterates over products in a Python loop. At $N{=}1000$ SKUs: \textbf{185 ms per query}.

\subsection{The Fix}
Full vectorized NumPy implementation of CIEDE2000. Pre-convert entire product DB to an $N \times 3$ Lab array offline. At runtime: \textbf{1.1 ms per query} ($168\times$ speedup).


\section{Bug \#4: PnP Without Camera Intrinsics $K$ (High)}

\subsection{The Error}
\texttt{cv2.solvePnP()} requires the intrinsic camera matrix:
\begin{equation}
K = \begin{bmatrix} f_x & 0 & c_x \\ 0 & f_y & c_y \\ 0 & 0 & 1 \end{bmatrix}
\end{equation}
Without $K$, 3D pose estimation is mathematically undefined. The original pipeline called PnP without providing $K$.

\subsection{Actual Resolution}
Still-photo pipeline uses 2D landmark ratios directly --- no 3D pose estimation needed. PnP is entirely removed from the pipeline.


\section{Bug \#5: Silent Dictionary Key Collision (Medium)}

Documented in Section 3 (Stage 6). Fixed with explicit \texttt{occ.keys() \& tex.keys()} collision detection before merge.


% ================================================================
% PART V: OPTICAL CORRECTIONS
% ================================================================
\part{Optical \& Color Science Corrections}

\section{Bradford Chromatic Adaptation Transform (D50 $\to$ D65)}
\label{sec:bradford}

\subsection{The Problem}
\begin{itemize}
    \item User selfie skin measurement: captured under ambient \textbf{D65} illuminant ($6500$K)
    \item Product catalog photography: shot under studio \textbf{D50} illuminant ($5000$K)
    \item Direct $\Delta E_{00}$ comparison without illuminant normalization: systematic error of $\Delta E_{00} \approx 3.5$--$6.0$
\end{itemize}

This means even a \textit{perfect} color-matching algorithm will recommend wrong products because the input data lives in different illuminant spaces.

\subsection{The Solution}
Apply the Bradford Chromatic Adaptation Matrix to map all product colors from D50 $\to$ D65 \textbf{offline, once}, before building the product DB:

\begin{equation}
\begin{bmatrix} X_{D65} \\ Y_{D65} \\ Z_{D65} \end{bmatrix} =
M_{\text{Brad}}^{-1} \cdot
\text{diag}\!\left(\frac{X_{D65,w}}{X_{D50,w}},\; \frac{Y_{D65,w}}{Y_{D50,w}},\; \frac{Z_{D65,w}}{Z_{D50,w}}\right) \cdot
M_{\text{Brad}} \cdot
\begin{bmatrix} X_{D50} \\ Y_{D50} \\ Z_{D50} \end{bmatrix}
\end{equation}

where:
\begin{equation}
M_{\text{Bradford}} = \begin{bmatrix}
0.8951 & 0.2664 & -0.1614 \\
-0.7502 & 1.7135 & 0.0367 \\
0.0389 & -0.0685 & 1.0296
\end{bmatrix}
\end{equation}

D65 white point: $(X_w, Y_w, Z_w) = (95.047, 100.0, 108.883)$\\
D50 white point: $(X_w, Y_w, Z_w) = (96.422, 100.0, 82.521)$

\begin{lstlisting}[language=Python, caption={Bradford D50 $\to$ D65 Adaptation}]
def bradford_d50_to_d65(xyz_d50):
    """Bradford Chromatic Adaptation: D50 -> D65."""
    M_bradford = np.array([
        [ 0.8951,  0.2664, -0.1614],
        [-0.7502,  1.7135,  0.0367],
        [ 0.0389, -0.0685,  1.0296]
    ])
    M_inv = np.linalg.inv(M_bradford)
    lms_scale = np.diag([95.047/96.422, 100.0/100.0, 108.883/82.521])
    M_cat = M_inv @ lms_scale @ M_bradford
    return (M_cat @ xyz_d50.T).T
\end{lstlisting}


\section{Vectorized CIEDE2000 ($\Delta E_{00}$) Product Matcher}

\subsection{What CIEDE2000 Is}
The most perceptually accurate color-difference formula standardized by CIE (2001). Unlike simple Euclidean distance in Lab space ($\Delta E_{ab}^*$), CIEDE2000 accounts for:
\begin{itemize}
    \item Lightness weighting ($S_L$)
    \item Chroma weighting ($S_C$)
    \item Hue weighting ($S_H$)
    \item Chroma-dependent hue rotation ($R_T$)
    \item The human eye's reduced sensitivity in the blue region
\end{itemize}

\subsection{Perceptual Thresholds}

\begin{table}[H]
\centering
\caption{$\Delta E_{00}$ Perceptual Interpretation for Cosmetics}
\begin{tabular}{@{}c l l@{}}
\toprule
$\Delta E_{00}$ & \textbf{Perception} & \textbf{Cosmetic Relevance} \\
\midrule
$< 1.0$ & Imperceptible & Perfect shade match \\
$1.0$--$2.0$ & Barely perceptible & Acceptable match for lips/blush \\
$2.0$--$3.5$ & Noticeable & May look ``off'' in certain lighting \\
$3.5$--$5.0$ & Obvious difference & Wrong shade; user will notice \\
$> 5.0$ & Very different & Completely wrong product \\
\bottomrule
\end{tabular}
\end{table}


\section{ROC/AUC Calibration Framework for 12 Seasons}
\label{sec:calibration}

\subsection{Why This Is Needed}
Current season classification uses arbitrary rule-based thresholds (e.g., ``if $a^* > 12$ and $b^* > 16$ then Warm''). These have never been validated against ground-truth expert annotations.

\subsection{Protocol}
\begin{enumerate}
    \item \textbf{Collect certified dataset}: $N \geq 600$ face images (50 per season), annotated under D65 lighting by certified Personal Color Analysis (PCA) consultants
    \item \textbf{Feature extraction}: Compute $(L^*, a^*, b^*, ITA^{\circ}, M, H)$ for each image
    \item \textbf{Threshold optimization}: Grid search over threshold vector $\mathbf{\Theta} = \{\theta_L, \theta_a, \theta_b, \theta_{ITA}, \theta_{Mel}, \theta_{Hb}\}$ maximizing Macro-F1:
    \begin{equation}
    \mathbf{\Theta}^* = \arg\max_{\mathbf{\Theta}} \frac{1}{12} \sum_{c=1}^{12} F1_c(\mathbf{\Theta})
    \end{equation}
    \item \textbf{Validation}: Plot per-class ROC curves and report AUC to establish defensible accuracy claims
\end{enumerate}

\begin{infoblock}[Phase 3 Deliverable]
This calibration framework is planned for Phase 3 of the roadmap. Phase 1 ships with rule-based heuristics and a prominent disclaimer that season accuracy is unvalidated.
\end{infoblock}


% ================================================================
% PART VI: STATUS & ROADMAP
% ================================================================
\part{Current Status \& Roadmap}

\section{What Exists Right Now}

\begin{table}[H]
\centering
\caption{Current Deliverables Inventory}
\begin{tabular}{@{}l c l@{}}
\toprule
\textbf{Artifact} & \textbf{Status} & \textbf{Notes} \\
\midrule
LaTeX Technical Report (IEEE format) & \checkmark & \texttt{makeup\_app\_technical\_report.tex/.pdf} \\
This Comprehensive Report & \checkmark & \texttt{glowup\_comprehensive\_report.tex/.pdf} \\
Product Database CSV (60 products) & \checkmark & \texttt{makeup\_app/data/product\_database.csv} \\
Bug Audit (5/5 resolved) & \checkmark & Documented in both reports \\
Bradford CAT code & \checkmark & In report listings (not yet in notebook) \\
Vectorized CIEDE2000 code & \checkmark & In report listings (not yet in notebook) \\
Strategy Pattern code & \checkmark & In report listings (not yet in notebook) \\
\midrule
Jupyter Notebook & $\times$ & \textbf{Not yet built} \\
ML Season Classifier & $\times$ & Phase 3 \\
Face Shape ML Classifier & $\times$ & Phase 3 \\
Labeled Training Dataset & $\times$ & Phase 3 \\
Gradio / Web Demo & $\times$ & Phase 4 \\
Production App & $\times$ & Phase 5 \\
\bottomrule
\end{tabular}
\end{table}


\section{5-Phase Roadmap}

\subsection{Phase 1: Working Prototype (1--2 weeks)}
Build the Jupyter Notebook with 9 cells covering the full pipeline. Use rule-based classifiers. Test on 10--20 selfies.\\
\textbf{Goal:} Photo in $\to$ Makeup card out.

\subsection{Phase 2: Product DB \& Visualization (1 week)}
Expand product DB to 200+ products. Add rich visualization: annotated face overlay, color palette cards, step-by-step visual guide.

\subsection{Phase 3: ML Classifier Upgrade (2--4 weeks)}
Collect/annotate labeled dataset. Train Random Forest or XGBoost for face shape and season. Validate with ROC/AUC. Replace rule-based thresholds.

\subsection{Phase 4: Interactive Demo (1 week)}
Wrap pipeline in Gradio or Streamlit. Upload photo $\to$ get results in browser. Shareable demo link.

\subsection{Phase 5: Production Deployment (future)}
Port to web app (React + Vite + MediaPipe WASM) or mobile (Flutter + TFLite). Separate project scope.


\section{Risk Assessment}

\begin{table}[H]
\centering
\caption{Key Risks and Mitigations}
\begin{tabular}{@{}p{5cm} c p{6cm}@{}}
\toprule
\textbf{Risk} & \textbf{Severity} & \textbf{Mitigation} \\
\midrule
12-season thresholds are wrong (no empirical validation) & High & Phase 3: collect labeled data, train ML classifier, validate with Macro-F1 \\
Lighting variation in selfies & Medium & Gray World white balance; instruct users to use natural daylight \\
Face shape rules too simple & Medium & Upgrade to ML classifier; test on diverse ethnicities \\
Product DB has incorrect hex values & Medium & Cross-reference multiple sources; visual swatch validation \\
Asian vs. Western beauty standards & Medium & Build demographic-aware modifier layer \\
Scope creep (AR try-on, video mode) & Medium & Defer to Phase 5 \\
\bottomrule
\end{tabular}
\end{table}


% ================================================================
% REFERENCES
% ================================================================
\begin{thebibliography}{99}
\bibitem{ita} P.~Chardon, I.~Cretois, and G.~Hourseau, ``Skin color typology and suntan measurement by Individual Typology Angle ($ITA^{\circ}$),'' \textit{International Journal of Cosmetic Science}, vol.~13, no.~4, pp.~191--208, 1991.
\bibitem{bradford} M.~D.~Fairchild, \textit{Color Appearance Models}, 3rd~ed. Chichester, UK: John Wiley \& Sons, 2013.
\bibitem{ciede2000} M.~R.~Luo, G.~Cui, and B.~Rigg, ``The development of the CIE 2000 colour-difference formula: CIEDE2000,'' \textit{Color Research \& Application}, vol.~26, no.~5, pp.~340--350, 2001.
\bibitem{mediapipe} Y.~Kartynnik, A.~Ablavatski, I.~Grishchenko, and M.~Grundmann, ``Real-time facial surface geometry estimation of single 2D images,'' \textit{arXiv:1907.06724}, 2019.
\bibitem{grayworld} G.~Buchsbaum, ``A spatial processor model for color constancy,'' \textit{Journal of the Franklin Institute}, vol.~310, no.~1, pp.~1--26, 1980.
\bibitem{srgb} IEC~61966-2-1:1999, ``Multimedia systems and equipment --- Colour measurement and management --- Part 2-1: Colour management --- Default RGB colour space --- sRGB.''
\bibitem{pca_seasons} C.~Jackson, \textit{Color Me Beautiful}, Ballantine Books, 1987.
\end{thebibliography}

\end{document}
